\documentclass[11pt,a4paper]{article}

\usepackage[margin=2.2cm]{geometry}
\usepackage{booktabs}
\usepackage{array}
\usepackage{enumitem}
\usepackage{titlesec}
\usepackage{xcolor}
\usepackage{fancyhdr}

\definecolor{entryrule}{RGB}{60,60,60}

% One heading per notebook entry: a horizontal rule, the date in bold,
% and the running notebook/page number on the right — the running-log
% shape an actual bench notebook is kept in, distinct from a report's
% \section structure.
\newcounter{entrynum}
\newcommand{\entry}[1]{%
  \stepcounter{entrynum}%
  \par\noindent\rule{\linewidth}{0.6pt}\\[2pt]
  {\bfseries\large #1} \hfill {\small Entry \theentrynum}\\[4pt]
}

\pagestyle{fancy}
\fancyhf{}
\lhead{Notebook --- Servo Controller Tuning}
\rhead{\thepage}
\renewcommand{\headrulewidth}{0.4pt}

\title{Laboratory Notebook \\ \large Servo Controller Tuning Project}
\author{R. Adeyemi}
\date{}

\begin{document}

\begin{titlepage}
\centering
\vspace*{3cm}
{\Huge\bfseries Laboratory Notebook\par}
\vspace{0.5cm}
{\Large Servo Controller Tuning Project\par}
\vspace{2cm}
{\large R. Adeyemi\par}
{\large Mechatronics Lab, Bench 4\par}
\vfill
{\large Notebook opened: 3 March\par}
\end{titlepage}

\entry{Monday, 3 March}

Set up the test rig: a 12V DC gear motor with a quadrature encoder
(500 counts/rev) driving a small flywheel, controlled through a PID
loop running on the bench microcontroller at 1\,kHz. Goal for this
notebook: tune the PID gains for a step-response settling time under
200\,ms with no more than 10\% overshoot.

Baseline run with $K_p = 1.0$, $K_i = 0$, $K_d = 0$: motor reaches
target speed but oscillates continuously, never settles. Expected ---
proportional-only control on a system with this much rotational
inertia always rings.

\begin{table}[h]
\centering
\caption{Step response, proportional-only gains}
\begin{tabular}{cccc}
\toprule
$K_p$ & Overshoot & Settling time & Notes \\
\midrule
0.5 & 61\% & did not settle & sustained oscillation \\
1.0 & 74\% & did not settle & sustained oscillation \\
2.0 & 88\% & did not settle & oscillation grows \\
\bottomrule
\end{tabular}
\end{table}

Higher $K_p$ makes it worse, not better --- as expected, need
derivative action to damp this.

\entry{Tuesday, 4 March}

Added derivative term. $K_p = 1.0$, $K_d = 0.05$: oscillation damps
out over about 6 cycles, settles at 340\,ms with 22\% overshoot.
Better, but overshoot is still too high and settling time is over the
200\,ms target.

Tried increasing $K_d$ further to see if that alone gets us there:

\begin{table}[h]
\centering
\caption{Step response, $K_p = 1.0$ fixed, varying $K_d$}
\begin{tabular}{ccc}
\toprule
$K_d$ & Overshoot & Settling time \\
\midrule
0.05 & 22\% & 340\,ms \\
0.10 & 9\%  & 265\,ms \\
0.15 & 4\%  & 310\,ms \\
\bottomrule
\end{tabular}
\end{table}

$K_d = 0.10$ hits the overshoot target (9\% $<$ 10\%) but settling
time is still 65\,ms over budget. Pushing $K_d$ further reduces
overshoot but slows the response down instead --- classic
overdamping. Need integral action to remove the small steady-state
droop I'm seeing at $K_d = 0.10$ before pushing $K_p$ up to speed
the response, otherwise a higher $K_p$ will just reopen the
overshoot problem $K_d$ just closed.

\entry{Wednesday, 5 March}

Added a small integral term to kill steady-state error, then raised
$K_p$ to speed up the response, re-tuning $K_d$ each time to hold
overshoot near the 10\% ceiling:

\begin{table}[h]
\centering
\caption{Full PID sweep}
\begin{tabular}{cccccc}
\toprule
$K_p$ & $K_i$ & $K_d$ & Overshoot & Settling time & Steady-state error \\
\midrule
1.0 & 0.02 & 0.10 & 8\%  & 258\,ms & $<$0.5\% \\
1.5 & 0.02 & 0.12 & 9\%  & 205\,ms & $<$0.5\% \\
2.0 & 0.02 & 0.15 & 10\% & 178\,ms & $<$0.5\% \\
2.5 & 0.02 & 0.18 & 15\% & 165\,ms & $<$0.5\% \\
\bottomrule
\end{tabular}
\end{table}

$K_p = 2.0$, $K_i = 0.02$, $K_d = 0.15$ is the best point found: 10\%
overshoot (right at the ceiling) and 178\,ms settling, comfortably
inside the 200\,ms target. $K_p = 2.5$ pushes settling time down
further but blows through the overshoot budget, so not usable.

\textbf{Chosen gains: $K_p = 2.0$, $K_i = 0.02$, $K_d = 0.15$.}

\entry{Thursday, 6 March}

Repeated the chosen-gain run 10 times to check reproducibility before
signing this off. Settling time varied 172--186\,ms, overshoot
9--11\% across runs --- consistent with the single-run result and
within acceptable trial-to-trial noise for this rig. One run at
11\% overshoot is a hair over the nominal 10\% ceiling but well
within the noise band observed here, not treated as a new data
point requiring re-tuning.

Next session: repeat this sweep at a higher flywheel load
(doubled moment of inertia) to check how much the chosen gains need
to change for a different mechanical load, since the production
application will not always be driving today's flywheel.

\end{document}
